\chapter{Securing WiFi: WPA2 and WPA3}

In Chapter~25, we saw how TLS secures a byte stream once TCP already connects a client to
a server. Before a single TCP packet reaches the router, however, your device has to
broadcast its frames across open air. Every time a laptop or phone wakes, sees a familiar
network name, and -- without anyone typing a password -- rejoins it and starts streaming
bytes, it runs a link-layer handshake first. That routine
reconnection uses one of the most widely deployed cryptographic protocols on the
planet, run billions of times a day, by coffee shops and airports and the cheap
router behind your television. It is also a protocol that has been broken, in
public and embarrassingly, more than once. WEP was badly broken. WPA2 stood for
thirteen years until KRACK forced a global patch. WPA3 is the current standard,
and it has had its own vulnerabilities.

This chapter is about the radio link: the single hop between your device and the
access point. That hop is special. Unlike a TCP connection across the Internet,
the WiFi link is broadcast. Every frame you send is, physically, audible to
anyone within range who cares to listen, whether or not they are ``on'' the
network. That changes the threat model in concrete ways, and it is why WiFi
encryption matters even when the site you are visiting already uses TLS.

We will build the WPA2 four-way handshake from the ground up -- where the keys
come from, what each of the four messages does, and why the design is shaped the
way it is. Then we will break it: WEP first (a clinic in how \emph{not} to use a
stream cipher), then KRACK against WPA2 (a clinic in how a single reused nonce can
unravel a careful protocol). Finally WPA3 and its Dragonfly/SAE handshake, which
replaces the weak link -- the password -- with a password-authenticated key
exchange. Throughout, we keep asking the practical question that actually decides
whether you should worry: what can an attacker on the same network, or just
within radio range, really do?

\chapterline

\section{The shape of the problem}

\begin{intuitionbox}
Alice has a laptop. Bob runs the access point (AP) -- the box that bridges the
WiFi radio to the wired Internet. Eve is sitting two tables over with a $\$30$
USB WiFi adapter in monitor mode, recording every frame in the air. The network
is named \texttt{CoffeeHouse} and is protected by a password that Alice typed in
once, months ago.

Three things must happen for the link to be secure against Eve:
\begin{enumerate}
  \item \textbf{Confidentiality.} Eve must not be able to read the contents of
        Alice's frames.
  \item \textbf{Integrity.} Eve must not be able to modify a frame -- flip a bit
        in a packet, redirect a DNS lookup -- without detection.
  \item \textbf{Authentication.} Alice must be talking to the real Bob, not to
        Eve impersonating the AP; and Bob must be talking to a device that knows
        the password.
\end{enumerate}
The list does \emph{not} include hiding \emph{that} Alice is using the
network, or hiding which AP she joins, or hiding how much she sends. WiFi headers
are sent in the clear. The link layer protects payloads, not metadata.
\end{intuitionbox}

The shared secret everyone starts from is the network password, which in WPA2
``Personal'' mode is called the \emph{passphrase} -- the thing on the sticker
under the router or written on the cafe chalkboard. In enterprise mode (WPA2-
Enterprise, used on corporate and university networks) the password is replaced
by an 802.1X/EAP exchange against a RADIUS server, and each user gets a fresh key
from that exchange. We focus on Personal mode, because it is where the
interesting cryptography lives and where the famous attacks land. The enterprise
key hierarchy is the same below the top; only the source of the master key
differs.

\begin{focusbox}
WiFi security uses many overlapping names, so here they are in one place.
\begin{itemize}
  \item \textbf{WEP} (1997): the original, fatally broken scheme. RC4 plus a
        24-bit IV.
  \item \textbf{WPA} (2003): a stopgap. Same RC4 hardware, but wrapped in TKIP
        to fix WEP's worst sins. Also broken now; treat as dead.
  \item \textbf{WPA2} (2004): the long-lived standard. AES in CCM mode (``CCMP'')
        for confidentiality and integrity, with the four-way handshake to derive
        keys. This is what ``WiFi encryption'' meant for over a decade.
  \item \textbf{WPA3} (2018): replaces the password-based handshake with SAE
        (``Dragonfly''), a password-authenticated key exchange, and adds
        forward secrecy. CCMP becomes GCMP-256 in the high-security tier.
  \item \textbf{IEEE 802.11i}: the standards document that defines WPA2's
        cryptography. ``RSN'' (Robust Security Network) is the same thing in the
        spec's language.
\end{itemize}
\end{focusbox}

\section{WEP, and how to break a stream cipher by reusing it}

Before WPA2 there was WEP, Wired Equivalent Privacy, and it is worth dissecting
because every mistake in it is a textbook lesson. WEP took RC4 -- a fine stream
cipher for its day -- and used it in just about the worst way one can.

Recall how a stream cipher works. It stretches a key into a long pseudorandom
keystream $Z$, and encrypts by XOR: $c = m \xor Z$. The basic rule of stream
ciphers is: \emph{never reuse a keystream}. If two
messages $m_1$ and $m_2$ are encrypted under the same keystream $Z$, then
\[
  c_1 \xor c_2 \;=\; (m_1 \xor Z) \xor (m_2 \xor Z) \;=\; m_1 \xor m_2,
\]
and the keystream cancels. The attacker now has the XOR of two plaintexts, which
-- with English text, known headers, or any structure -- is often enough to
recover both.

So the entire security of a stream-cipher system reduces to: \emph{is the
keystream ever reused?} WEP guaranteed that it would be.

\begin{whybox}
WEP's per-frame key was $K = \text{IV} \concat k$, where $k$ is the long-term
shared WEP key (the password, essentially) and IV is a 24-bit initialization
vector sent in the clear with each frame. RC4 was keyed with $K$ to produce the
keystream for that frame.

The IV is 24 bits, so there are only $2^{24} \approx 16.7$ million possible
values. On a busy network sending, say, $1000$ frames per second, you exhaust the
entire IV space in about $2^{24}/1000 \approx 16{,}777$ seconds, under five
hours. After that, IVs \emph{must} repeat, and every repeat is a keystream reuse,
because the long-term $k$ never changes. By the birthday bound (see the companion
probability volume), collisions start appearing far sooner -- after roughly
$\sqrt{2^{24}} = 2^{12} = 4096$ frames you expect your first repeated IV, which is a few
seconds of traffic.
\end{whybox}

\begin{attackbox}
\textbf{WEP is broken three ways over, and any one would be fatal.}

\emph{Keystream reuse.} As above, repeated IVs give $m_1 \xor m_2$. Many WiFi
frames have predictable structure (ARP packets are a fixed size with known
fields; IP headers are guessable), so the attacker often knows or can guess
$m_1$, recovering $m_2$ outright -- and recovering the keystream for that IV,
which then decrypts every future frame using the same IV.

\emph{No real integrity.} WEP ``protects'' integrity with a CRC-32 checksum,
which is \emph{linear}: $\mathrm{CRC}(x \xor \delta) = \mathrm{CRC}(x) \xor
\mathrm{CRC}(\delta)$. Because encryption is also linear (XOR), an attacker can
flip any bit of the ciphertext and compute exactly how to fix the checksum so the
frame still verifies. CRC is an error-detecting code designed for random
errors, not a MAC, so it gives no protection against a deliberate attacker.

\emph{The RC4 key-schedule weakness (FMS).} In 2001, Fluhrer, Mantin, and Shamir
showed that RC4's key schedule leaks information about the key when the first few
keystream bytes are known and the key has a particular structure -- exactly the
structure WEP's $\text{IV} \concat k$ provides, with the attacker-visible IV as
the variable prefix. Each captured frame leaks a tiny bias about $k$. Collect
enough ``weak IV'' frames and you solve for $k$ directly. Later tools (KoreK,
then the PTW attack of 2007) drove the requirement down to tens of thousands of
frames -- a minute or two of traffic, less if you inject packets to force more.
The recovered key is the \emph{actual} WEP key, decrypting everything past and
future.
\end{attackbox}

RC4's biases were a real flaw, but WEP would
have fallen to keystream reuse and the CRC malleability even with a perfect
cipher. A cipher is only as good as the construction around it.
WEP got the IV size wrong, got the integrity primitive wrong (a checksum where a
MAC was needed), and rekeyed never. WPA2 fixes all three: a long never-repeating
counter, a real MAC, and fresh session keys per association.

\begin{historybox}
WEP shipped in 1997 inside the original 802.11 standard, designed by a committee
optimizing for cheap hardware and an export regime that capped key sizes. The 24-
bit IV and 40-bit key were partly artifacts of 1990s U.S. crypto export limits.
By 2001 the FMS paper was public; by the mid-2000s, point-and-click tools
(\texttt{aircrack-ng}) cracked WEP in minutes. The 2007 breach of TJX (millions
of credit cards) was traced to a store WiFi network still running WEP. The
standard formally deprecated it, and modern devices have removed WEP support
entirely. If you ever see WEP offered, treat the network as plaintext.
\end{historybox}

\section{WPA2: the key hierarchy}

WPA2 starts from the same place -- a shared passphrase -- but never uses it
directly to encrypt anything. Instead it derives a hierarchy of keys, with the
passphrase at the root and short-lived per-frame protection at the leaves. The
guiding principle: the long-term secret should touch the air as little as
possible, and every session should use independent keys so that breaking one
tells you nothing about another.

\subsection{From passphrase to PMK}

The passphrase is first stretched into a 256-bit \emph{Pairwise Master Key}
(PMK). In WPA2-Personal the stretching function is PBKDF2 (Password-Based Key
Derivation Function 2) with HMAC-SHA1, $4096$ iterations, salted by the network
name (SSID):
\[
  \text{PMK} \;=\; \mathrm{PBKDF2}\bigl(\text{passphrase},\ \text{SSID},\
  4096,\ 256\ \text{bits}\bigr).
\]

\begin{whybox}
Two design choices in that line matter.

The \textbf{4096 iterations} make each password guess expensive. An attacker who
captures a handshake and wants to brute-force the passphrase must run PBKDF2 --
$4096$ HMAC-SHA1 calls -- per candidate. Without it, a GPU could test billions of
guesses per second; with it, the rate drops by three to four orders of magnitude.
This is the same key-stretching idea you saw for password hashing. It does not
make a weak passphrase safe, only slower to crack.

The \textbf{SSID as salt} means the same passphrase on two networks with
different names yields different PMKs. This blocks precomputation across networks
-- almost. The catch is that the salt is the \emph{network name}, not a random
value, and popular names (\texttt{linksys}, \texttt{xfinitywifi}, the default
names of common routers) are predictable. Attackers precompute giant tables
(``rainbow tables'') of PMKs for common SSID-plus-passphrase pairs. If your SSID
is a default and your passphrase is in a wordlist, the PMK is effectively
already computed.
\end{whybox}

\begin{connectionbox}
PBKDF2's cost is exactly the modular-exponentiation-versus-hashing tradeoff the
algorithms volume analyzes for password storage: you deliberately pick a slow
function so the legitimate party pays once at login but the attacker pays per
guess. The birthday-bound reasoning behind WEP's IV collisions, and behind the
nonce-reuse risks we will meet in GCMP, is the same $\sqrt{N}$ collision law
derived in the probability volume's chapter on the birthday problem.
\end{connectionbox}

\subsection{From PMK to PTK: the four-way handshake}

The PMK is a long-term per-network secret -- everyone who knows the passphrase can
compute the same PMK. (Note that in WPA2-Personal, your neighbor on the same
cafe network has the \emph{same PMK as you}. We will return to what that means for
attacks.) The PMK is never used to encrypt traffic directly. Instead, each time a
device associates, a fresh \emph{Pairwise Transient Key} (PTK) is derived for that
session, mixing in random nonces from both sides so that every association gets
independent keys.

The PTK is computed by a PRF (in WPA2, an HMAC-SHA1-based key-derivation
function) from five inputs:
\[
  \text{PTK} \;=\; \mathrm{PRF}\bigl(\text{PMK};\ \text{Anonce} \concat
  \text{Snonce} \concat \mathrm{AA} \concat \mathrm{SA}\bigr),
\]
where Anonce is a random nonce chosen by the AP (the authenticator), Snonce is a
random nonce chosen by the station (your device), AA is the AP's MAC address, and
SA is the station's MAC address. The two MAC addresses are sorted into a fixed
order, as are the two nonces, so both sides feed the PRF identical input and get
identical output.

The PTK is then sliced into several sub-keys:
\begin{itemize}
  \item \textbf{KCK} (Key Confirmation Key, 128 bits): used to MAC the handshake
        messages themselves, proving each side derived the same PTK.
  \item \textbf{KEK} (Key Encryption Key, 128 bits): used to encrypt the group
        key when it is handed to the station.
  \item \textbf{TK} (Temporal Key, 128 bits): the actual key that protects data
        frames with AES-CCMP.
\end{itemize}

\begin{intuitionbox}
Why nonces from \emph{both} sides? If only the AP contributed a nonce, then two
sessions where the AP happened to reuse a nonce would derive the same PTK -- and a
device that always sends the same Snonce would help that along. Mixing two
independent random values means the PTK is fresh as long as \emph{either} side
picks a fresh nonce, and the MAC addresses bind the key to this specific pair of
devices. The handshake is a Diffie-Hellman in spirit -- both parties contribute
randomness -- except the shared secret is the pre-existing PMK rather than a fresh
DH exchange. That distinction causes WPA2's main weakness: there is no
forward secrecy, because the PMK is fixed. Compromise the passphrase later and you
can recompute every PTK from recorded nonces, decrypting traffic you captured
years ago. WPA3 fixes exactly this.
\end{intuitionbox}

Figure~\ref{ch26:fig-wpa2-handshake} details the four-way handshake message ladder.
It is called ``four-way'' because four EAPOL (Extensible Authentication Protocol over LAN)
frames cross the link.

\begin{figure}[H]
\centering
\begin{tikzpicture}[font=\scriptsize, >=stealth]
  % Endpoints
  \node[draw=aliceblue!80, fill=aliceblue!15, rounded corners=2pt, minimum width=2.8cm, minimum height=0.7cm, align=center] (sta) at (0,0) {
    \textbf{Station (Alice)}\\
    {\tiny Holds PMK $= \operatorname{PBKDF2}(\pi, \text{SSID})$}
  };
  \node[draw=bobgreen!80, fill=bobgreen!15, rounded corners=2pt, minimum width=2.8cm, minimum height=0.7cm, align=center] (ap) at (8.5,0) {
    \textbf{Access Point (Bob)}\\
    {\tiny Holds PMK $= \operatorname{PBKDF2}(\pi, \text{SSID})$}
  };

  \draw[thick, gray!60] (sta) -- (0, -7.0);
  \draw[thick, gray!60] (ap) -- (8.5, -7.0);

  % M1: Anonce
  \draw[->, thick, bobgreen!80!black] (8.5, -1.0) -- node[above, font=\tiny]{1. \textbf{M1}: Anonce (unprotected)} (0, -1.0);
  \node[draw=aliceblue!70, fill=aliceblue!10, rounded corners=1pt, font=\tiny, align=left, text width=3.8cm] at (0, -1.8) {
    Pick Snonce.\\
    Derive $\text{PTK} = \operatorname{PRF}(\text{PMK},$\\
    $\quad\text{Anonce}, \text{Snonce}, \dots)$\\
    Split PTK: KCK, KEK, TK.
  };

  % M2: Snonce + MIC
  \draw[->, thick, aliceblue!80!black] (0, -2.8) -- node[above, font=\tiny]{2. \textbf{M2}: Snonce + $\operatorname{MIC}_{\mathrm{KCK}}(\text{M2})$} (8.5, -2.8);
  \node[draw=bobgreen!70, fill=bobgreen!10, rounded corners=1pt, font=\tiny, align=left, text width=3.8cm] at (8.5, -3.6) {
    Derive PTK from Anonce, Snonce.\\
    Verify $\operatorname{MIC}_{\mathrm{KCK}}(\text{M2})$.\\
    $\Longrightarrow$ Alice knows passphrase.
  };

  % M3: GTK + MIC
  \draw[->, thick, bobgreen!80!black] (8.5, -4.5) -- node[above, font=\tiny]{3. \textbf{M3}: $\Enc_{\mathrm{KEK}}(\text{GTK}) + \operatorname{MIC}_{\mathrm{KCK}}(\text{M3})$} (0, -4.5);
  \node[draw=aliceblue!70, fill=aliceblue!10, rounded corners=1pt, font=\tiny, align=left, text width=3.8cm] at (0, -5.2) {
    Verify $\operatorname{MIC}_{\mathrm{KCK}}(\text{M3})$.\\
    Install PTK (TK) and GTK.\\
    Initialize packet number $\text{PN} = 0$.
  };

  % M4: Ack + MIC
  \draw[->, thick, aliceblue!80!black] (0, -6.0) -- node[above, font=\tiny]{4. \textbf{M4}: Ack + $\operatorname{MIC}_{\mathrm{KCK}}(\text{M4})$} (8.5, -6.0);
  \node[draw=bobgreen!70, fill=bobgreen!10, rounded corners=1pt, font=\tiny, align=left, text width=3.8cm] at (8.5, -6.6) {
    Verify MIC. Install PTK.
  };

  % Secure data flow
  \draw[<->, thick, structurepurple] (0, -7.3) -- node[above, font=\tiny\bfseries]{Encrypted WiFi Data: AES-CCMP under Temporal Key (TK)} (8.5, -7.3);
\end{tikzpicture}
\caption{The WPA2 four-way handshake. Alice and Bob mutually authenticate their possession of the
Pairwise Master Key (PMK) by exchanging nonces, deriving the Pairwise Transient Key (PTK),
and verifying Message Integrity Codes (MICs) before installing the data encryption key (TK).}
\label{ch26:fig-wpa2-handshake}
\end{figure}

\begin{enumerate}
  \item[\textbf{M1}] The AP sends its random \textbf{Anonce} to the station. This
        message is unauthenticated -- the station cannot yet tell a real AP from a
        forgery, because no key exists yet. (Remember this for KRACK.)
  \item[\textbf{M2}] The station picks its \textbf{Snonce}, now has both nonces
        and both MACs, and computes the PTK. It sends Snonce back, with a
        \textbf{MIC} (Message Integrity Code) computed under the KCK. The AP can
        now also compute the PTK; if the MIC verifies, the AP knows the station
        derived the same PTK, hence knows the PMK, hence knows the passphrase.
  \item[\textbf{M3}] The AP sends the \textbf{GTK} (Group Temporal Key, the key
        for broadcast/multicast frames), encrypted under the KEK, with its own
        MIC under the KCK. The station verifies the MIC: now \emph{it} knows the
        AP knew the PMK. Mutual authentication is complete. The station installs
        the PTK for use.
  \item[\textbf{M4}] The station sends an acknowledgment (with a MIC). On receipt
        the AP installs the PTK. Both sides switch to encrypted data frames.
\end{enumerate}

\begin{focusbox}
The four-way handshake authenticates by \emph{proof of knowledge of the PMK},
never by sending the PMK or passphrase. Each side proves it can compute the
correct MIC under a key (KCK) that only a holder of the PMK could derive, given
the public nonces. This is a challenge-response: the nonces are the challenges,
the MICs are the responses. Nothing secret crosses the air -- which is exactly why
an eavesdropper who records the whole handshake \emph{still} cannot read the
traffic, unless they also know the passphrase.
\end{focusbox}

\subsection{Protecting data: AES-CCMP}

Once the TK is installed, each data frame is protected with CCMP: AES in
\textbf{Counter mode} for confidentiality, with \textbf{CBC-MAC} for integrity --
the ``CCM'' authenticated-encryption mode. This is the piece WEP got so wrong, and
WPA2 gets right.

\begin{whybox}
CCMP solves WEP's three failures cleanly:
\begin{itemize}
  \item \textbf{No keystream reuse.} The AES-CTR counter includes a 48-bit
        \emph{packet number} (PN) that increments on every frame and never
        repeats for the life of the TK. $2^{48}$ frames is astronomically more
        than any session will send, and a fresh TK each association resets it. The
        PN is also checked for monotonic increase, blocking replay.
  \item \textbf{Real integrity.} CBC-MAC under AES is a genuine MAC, not a linear
        checksum. An attacker who flips a ciphertext bit cannot repair the MIC
        without the key. Forgery probability is roughly $2^{-64}$ per try.
  \item \textbf{Authenticated headers.} The MIC covers not just the payload but
        the frame header fields that matter (addresses, QoS), so an attacker
        cannot redirect a frame by editing its header.
\end{itemize}
The PN fixes the problem with WEP's 24-bit IV: it is long enough never to
repeat, included in the authenticated data, and checked for replay.
\end{whybox}

\begin{connectionbox}
CCM is the counter-plus-CBC-MAC authenticated encryption mode studied in
Chapter~16. Counter mode turns the block cipher AES into a stream cipher (so
the no-reuse rule applies, hence the careful PN), while CBC-MAC supplies the
integrity that CTR mode alone lacks. WPA3's high tier swaps CCM for GCM (GCMP-256),
the same AES-GCM you meet in TLS 1.3 -- faster, but with an even sharper nonce-
reuse cliff, which is why the 256-bit version and careful PN handling matter.
\end{connectionbox}

\section{A worked miniature handshake}

Let us trace the key derivation with tiny stand-in numbers, so the data flow is
concrete. Real WPA2 uses 256-bit PMKs and HMAC-SHA1; we will use a toy hash so
you can follow the arithmetic. This is a \emph{model}, not the real algorithm --
the point is the shape of the dependency graph, not the bytes.

\begin{example}
\label{ch26:toy-handshake}
Let our toy ``PRF'' be $\mathrm{prf}(K, x) = (K \cdot x + 7) \bmod 101$, working
mod the prime $101$. Suppose:
\begin{itemize}
  \item Passphrase-derived $\text{PMK} = 42$ (in reality, $256$ bits from PBKDF2).
  \item AP nonce $\text{Anonce} = 13$, station nonce $\text{Snonce} = 28$.
  \item We combine nonces by addition for the toy: $x = \text{Anonce} +
        \text{Snonce} = 41$.
\end{itemize}
Then the toy PTK is
\[
  \text{PTK} = \mathrm{prf}(42, 41) = (42 \cdot 41 + 7) \bmod 101
             = (1722 + 7) \bmod 101 = 1729 \bmod 101.
\]
Now $1729 = 17 \cdot 101 + 12$, so $\text{PTK} = 12$.

Slice it (toy-style) into a confirmation key and a data key: say
$\text{KCK} = \text{PTK} \bmod 10 = 2$ and $\text{TK} = \lfloor \text{PTK}/10
\rfloor = 1$ together with the nonces. The station's M2 MIC is
$\mathrm{prf}(\text{KCK}, \text{Snonce}) = (2 \cdot 28 + 7) \bmod 101 = 63$. The AP
recomputes the same PTK (it has PMK and both nonces), derives the same KCK $= 2$,
recomputes $\mathrm{prf}(2, 28) = 63$, and sees the MIC matches -- so the station
proved knowledge of the PMK without revealing it.
\end{example}

Trace what an eavesdropper sees: Anonce $= 13$, Snonce $= 28$, the MIC $= 63$, the
MAC addresses. To forge M2 or read traffic, Eve needs the PTK, which needs the
PMK $= 42$, which she does not have. Her only avenue is to guess the passphrase
offline: pick a candidate, run PBKDF2 to a candidate PMK, derive a candidate PTK
and KCK, and check whether $\mathrm{prf}(\text{KCK}, 28)$ equals the captured MIC
$63$. That check is the entire WPA2-Personal offline attack, and the next section
shows how it works in practice.

\begin{trapbox}
Do not read Example~\ref{ch26:toy-handshake} as the real PRF. Real WPA2 uses
HMAC-SHA1 over a precisely ordered concatenation of sorted MAC addresses and
sorted nonces, expanded to produce $384$ or $512$ bits of PTK. The toy linear
function $K x + 7$ is invertible and would let an attacker solve for $K$ from one
(input, output) pair -- a catastrophe a real PRF prevents. The toy is for tracing
data flow only.
\end{trapbox}

\section{The offline dictionary attack on WPA2-Personal}

WPA2-Personal's authentication reduces to: ``prove you know the PMK by
producing the right MIC.'' Everything an attacker needs to \emph{check} a
passphrase guess is in the four-way handshake, which is sent in the clear. So
WPA2-Personal is vulnerable to an \emph{offline} dictionary attack, and there is
no fixing it within the protocol -- it is inherent to deriving keys from a low-
entropy password without a PAKE. (WPA3 is precisely the fix.)

\begin{attackbox}
\textbf{Capture-and-crack against WPA2-Personal.}

\emph{Step 1: capture a handshake.} Eve puts her adapter in monitor mode and
waits for any device to associate, recording M1 through M4. She does not need to
wait long -- she can \emph{force} it. A WiFi \textbf{deauthentication} frame is a
management frame that, in plain WPA2, is unauthenticated: the AP can tell a client
``you are disconnected,'' and clients obey. Eve spoofs a deauth from the AP's MAC
to Alice. Alice drops off and immediately reconnects, performing a fresh four-way
handshake that Eve records.

\emph{Step 2: crack offline.} From the handshake Eve has Anonce, Snonce, both
MACs, and the M2 MIC. For each candidate passphrase $p$ in a wordlist:
\begin{align*}
  \text{PMK} &\gets \mathrm{PBKDF2}(p,\ \text{SSID},\ 4096),\quad
  \text{PTK} \gets \mathrm{PRF}(\text{PMK}, \ldots), \\
  \text{check } &\mathrm{MIC}(\text{KCK}, \text{M2}) \stackrel{?}{=} \text{captured}.
\end{align*}
A match means $p$ is the passphrase. This is entirely offline -- no further
contact with the network, no rate limiting, as many guesses as Eve's hardware
allows.

\emph{What decides success.} The PBKDF2 cost ($4096$ iterations) is the only
brake, and modern GPUs push WPA2 candidate rates into the hundreds of thousands
to millions per second per card. A passphrase from a wordlist, a common pattern
(\texttt{Summer2024!}), a phone number, or anything under roughly $10$--$12$
random characters of entropy is realistically crackable. A long random passphrase
(say $16$+ random characters, or a five-word diceware phrase) is not. Security here
depends entirely on the entropy of the human-chosen password.
\end{attackbox}

\begin{focusbox}
This is the single most important practical fact about WPA2-Personal: \textbf{the
encryption is only as strong as the passphrase}, because anyone within radio range
can capture a handshake and grind it offline forever. AES-128 contributes nothing
if the passphrase is \texttt{password123}. Use long, random passphrases, and
prefer WPA3 (which defeats this attack entirely) where available.
\end{focusbox}

\begin{implementbox}
A learning-only sketch of the WPA2-Personal offline check. It computes the PMK and
the MIC over a captured handshake and tests passphrases from a wordlist. It is
deliberately incomplete and slow; real tools like \texttt{hashcat} do this on GPUs.
\textbf{Use this only to understand the attack on networks you own. Cracking
others' WiFi is illegal in most jurisdictions.}

\begin{lstlisting}[language=Python]
import hashlib, hmac

def pmk(passphrase, ssid):
    # PBKDF2-HMAC-SHA1, 4096 iters, 32-byte PMK -- the WPA2-Personal stretch.
    return hashlib.pbkdf2_hmac("sha1", passphrase.encode(),
                               ssid.encode(), 4096, 32)

def prf512(pmk, A, B):
    # 802.11i PRF: HMAC-SHA1 expansion to 64 bytes. "A" is a label string,
    # "B" = min(macA,macB)||max(macA,macB)||min(n1,n2)||max(n1,n2).
    out = b""
    i = 0
    while len(out) < 64:
        out += hmac.new(pmk, A + b"\x00" + B + bytes([i]),
                        hashlib.sha1).digest()
        i += 1
    return out[:64]  # PTK = KCK(16) || KEK(16) || TK(16) || extra

def ptk(pmk_bytes, ap_mac, sta_mac, anonce, snonce):
    B = min(ap_mac, sta_mac) + max(ap_mac, sta_mac) \
        + min(anonce, snonce) + max(anonce, snonce)
    return prf512(pmk_bytes, b"Pairwise key expansion", B)

def mic_matches(passphrase, ssid, ap_mac, sta_mac,
                anonce, snonce, eapol_m2, captured_mic):
    p = ptk(pmk(passphrase, ssid), ap_mac, sta_mac, anonce, snonce)
    kck = p[:16]
    # WPA2 (CCMP) uses HMAC-SHA1 for the MIC, truncated to 16 bytes.
    mic = hmac.new(kck, eapol_m2, hashlib.sha1).digest()[:16]
    return mic == captured_mic

# for guess in open("wordlist.txt"):
#     if mic_matches(guess.strip(), ssid, ap_mac, sta_mac,
#                    anonce, snonce, eapol_m2, captured_mic):
#         print("passphrase:", guess); break
\end{lstlisting}
The real attack also handles the EAPOL frame's MIC field being zeroed before the
computation, and the choice of HMAC-SHA1 versus AES-CMAC depending on the cipher
suite. Apart from those details, the check is the one shown above.
\end{implementbox}

\section{KRACK: breaking WPA2 without the password}

For thirteen years, the four-way handshake's cryptography held: if your passphrase
was strong, your traffic was safe. Then in 2017 Mathy Vanhoef published KRACK --
the Key Reinstallation Attack -- which broke WPA2 \emph{without knowing the
passphrase at all}. KRACK is notable because the flaw was
not in the cipher or the key derivation. It was in the \emph{state machine} --
specifically, in what happens when a handshake message is retransmitted.

\begin{intuitionbox}
Recall message M3: the AP sends the GTK and a MIC, and on receiving it the station
\emph{installs the PTK}, which resets the CCMP packet number (the nonce counter)
to its starting value and begins encrypting data frames. Message M4 acknowledges
this.

But networks are lossy. M4 might not arrive. The standard says: if the AP does not
get M4, it \emph{retransmits M3}. And the standard says the station, on receiving
a retransmitted M3, should \emph{reinstall} the same PTK -- and in doing so,
\textbf{reset the packet number counter back to its starting value.}

So reinstalling the key resets the nonce, and resetting the nonce in a
counter-mode cipher reuses keystream. KRACK forces
exactly this reset, on demand, by replaying M3.
\end{intuitionbox}

\begin{attackbox}
\textbf{KRACK against the four-way handshake.}

Eve sits as a man-in-the-middle at the radio layer, relaying frames between Alice
and the AP (she can do this with a cloned AP on a different channel -- a ``channel-
based MitM'' -- without knowing any key, because management and handshake frames
are unprotected).

\begin{enumerate}
  \item The handshake proceeds normally up to M3. Eve \emph{blocks} M4 from
        reaching the AP (she is in the middle).
  \item The AP, not hearing M4, retransmits M3. Eve forwards this retransmitted M3
        to Alice.
  \item Alice, following the standard, \textbf{reinstalls} the PTK and \textbf{
        resets her packet number to its initial value} -- even though she has
        already sent data frames with that key.
  \item Now Alice encrypts her next data frame with the \emph{same} PN (nonce) she
        used right after the first install. Two plaintexts, one keystream:
        $c_1 \xor c_2 = m_1 \xor m_2$. The keystream is recovered for those PNs,
        and with known plaintext (TCP/IP headers, retransmissions Eve can force)
        Eve decrypts the frames.
\end{enumerate}

On some implementations the damage was worse than keystream reuse. Android $6.0$
and certain \texttt{wpa\_supplicant} versions, on reinstall, set the key to
\emph{all zeros}, so Eve did not even need the password-derived keystream -- the
traffic was encrypted under a known, all-zero key. That turned KRACK from
``decrypt some frames'' into ``decrypt all captured frames without searching for the passphrase.''
\end{attackbox}

Figure~\ref{ch26:fig-krack} illustrates the channel-based man-in-the-middle sequence that triggers
this nonce reset and keystream reuse.

\begin{figure}[H]
\centering
\begin{tikzpicture}[font=\scriptsize, >=stealth]
  % Parties
  \node[draw=aliceblue!80, fill=aliceblue!15, rounded corners=2pt, minimum width=2.4cm, minimum height=0.7cm, align=center] (sta) at (0,0) {
    \textbf{Station (Alice)}\\
    {\tiny CCMP Nonce Counter: PN}
  };
  \node[draw=trapred!80, fill=trapred!15, rounded corners=2pt, minimum width=2.4cm, minimum height=0.7cm, align=center] (mitm) at (5.0,0) {
    \textbf{Attacker (Eve)}\\
    {\tiny Channel-based MitM}
  };
  \node[draw=bobgreen!80, fill=bobgreen!15, rounded corners=2pt, minimum width=2.4cm, minimum height=0.7cm, align=center] (ap) at (10.0,0) {
    \textbf{Access Point (Bob)}\\
    {\tiny Standard 802.11 AP}
  };

  \draw[thick, gray!60] (sta) -- (0, -6.5);
  \draw[thick, gray!60] (mitm) -- (5.0, -6.5);
  \draw[thick, gray!60] (ap) -- (10.0, -6.5);

  % Initial M3 delivery
  \draw[->, thick, bobgreen!80!black] (10.0, -0.9) -- node[above, font=\tiny]{1. M3 (GTK, MIC)} (5.0, -0.9);
  \draw[->, thick, trapred!80!black] (5.0, -0.9) -- (0, -0.9);
  \node[draw=aliceblue!70, fill=aliceblue!10, rounded corners=1pt, font=\tiny, align=left] at (0, -1.5) {
    Install PTK. Initialize $\text{PN} = 0$.
  };

  % Alice transmits M4, Eve BLOCKS it
  \draw[->, thick, aliceblue!80!black] (0, -2.1) -- node[above, font=\tiny]{2. M4 (Ack)} (5.0, -2.1);
  \node[draw=trapred!80, fill=trapred!20, rounded corners=1pt, font=\tiny\bfseries, text=trapred!90!black] at (5.0, -2.5) {$\times$ Eve BLOCKS M4};

  % Alice sends normal frame 1
  \draw[->, thick, structurepurple] (0, -2.9) -- node[above, font=\tiny]{Data Frame 1: $\Enc_{\mathrm{TK}}(\text{PN}=0, m_1)$} (5.0, -2.9);
  \draw[->, thick, structurepurple] (5.0, -2.9) -- (10.0, -2.9);
  \node[font=\tiny, gray!70] at (0, -3.3) {PN advances: $\text{PN} = 1, 2, \dots$};

  % Bob retransmits M3
  \draw[->, thick, bobgreen!80!black] (10.0, -3.8) -- node[above, font=\tiny]{3. AP retransmits M3 (no M4 ack)} (5.0, -3.8);
  \draw[->, thick, trapred!80!black] (5.0, -3.8) -- node[above, font=\tiny]{Forward retransmitted M3} (0, -3.8);

  % Key reinstallation bug
  \node[draw=trapred!80, fill=trapred!15, rounded corners=2pt, font=\tiny, align=left, text width=4.5cm] at (0, -4.7) {
    \textbf{Key Reinstallation Flaw}:\\
    Standard specifies reinstalling PTK on duplicate M3.\\
    $\Longrightarrow$ Station \textbf{resets $\text{PN} \leftarrow 0$}!
  };

  % Reused keystream
  \draw[->, thick, trapred!90!black] (0, -5.6) -- node[above, font=\tiny\bfseries, text=trapred!90!black]{4. Data Frame 2: $\Enc_{\mathrm{TK}}(\mathbf{PN=0}, m_2)$} (5.0, -5.6);

  % Keystream cancellation banner
  \node[draw=trapred!70, fill=trapred!10, rounded corners=2pt, font=\tiny, align=center, text width=11.5cm] at (5.0, -6.8) {
    \textbf{Catastrophic Keystream Reuse}: Both Frame 1 and Frame 2 encrypted with identical key and $\text{PN} = 0$.\\
    Eve computes $c_1 \xor c_2 = m_1 \xor m_2$. Using known header fields or web traffic structure,\\
    the keystream is stripped and plaintexts are recovered without ever cracking the WiFi password.
  };
\end{tikzpicture}
\caption{The KRACK (Key Reinstallation Attack) protocol flaw. By blocking Alice's Message 4 acknowledgment,
the attacker induces the Access Point to retransmit Message 3. Compliant 802.11 clients reinstall the
temporal key and reset the packet number counter ($\text{PN} \leftarrow 0$), causing catastrophic nonce reuse in CCMP/GCMP.}
\label{ch26:fig-krack}
\end{figure}

\begin{whybox}
Why did such a careful protocol have this hole? Because the cryptographic
\emph{messages} were authenticated (the MICs are correct, the PTK is right -- Eve
never forges anything), but the \emph{state machine around them} did not treat
``reinstall the same key'' as the dangerous operation it is. The protocol proof for
the four-way handshake, done years earlier, modeled the key agreement -- and it was
correct! The flaw lived in a layer the proof did not cover: the retransmission and
reinstall behavior mandated by the standard's prose, not its key-exchange core.

The fix is small and was deployed via software patches in late 2017: do not reset
the packet number on a key reinstall, or refuse to reinstall an already-installed
key. No new key material, no protocol redesign -- just ``never let a nonce go
backwards.'' KRACK did not retire WPA2; a patched WPA2 stack is not vulnerable.
But it forced an Internet-wide update and showed how completely counter-mode
AEAD depends on nonce discipline.
\end{whybox}

\begin{connectionbox}
KRACK is the WiFi instance of a general law you will see again in TLS, IPsec, and
QUIC: \emph{nonce reuse in an AEAD mode is catastrophic}. Chapter~16
shows that reusing a GCM nonce leaks not only plaintext (via keystream canceling)
but the authentication key itself, enabling forgeries. CCM is gentler -- nonce
reuse there ``only'' leaks plaintext, not the MAC key -- but the lesson is
identical: the counter must never repeat or reverse. KRACK is what happens when a
state machine quietly violates that invariant.
\end{connectionbox}

\section{WPA3 and the SAE/Dragonfly handshake}

WPA2-Personal has one structural limitation that no patch can fix: the offline
dictionary attack. Because authentication consists solely of proving knowledge of the
PMK, and the PMK comes deterministically from the passphrase, anyone who captures
a handshake can test candidate passphrases offline forever. That weakness is built
into deriving keys directly from a password.

WPA3-Personal closes this avenue by replacing the four-way handshake's front end
with a \textbf{Password-Authenticated Key Exchange} (PAKE), specifically
\textbf{SAE} -- Simultaneous Authentication of Equals, also called Dragonfly. A
PAKE lets two parties who share a (possibly weak) password agree on a strong,
random shared key such that:
\begin{itemize}
  \item An eavesdropper learns \emph{nothing} that helps an offline guess.
  \item An active attacker gets at most \emph{one} password guess per live
        interaction -- they must talk to the AP for each guess, and the AP can
        rate-limit. Offline grinding is impossible.
  \item The resulting session key has \textbf{forward secrecy}: it is a fresh
        Diffie-Hellman secret, so compromising the password \emph{later} does not
        decrypt traffic recorded earlier.
\end{itemize}

\begin{intuitionbox}
A PAKE folds the password into a Diffie-Hellman exchange so that
the DH messages look like random group elements to anyone who does not know the
password, and so that the exchange only completes if both sides used the same
password. You get the best of both: the strong randomness and forward secrecy of
Diffie-Hellman, plus authentication from the shared password -- without ever
sending anything an eavesdropper can test offline. The phrase ``of equals'' means
there is no client/server asymmetry: AP and station run identical roles, either
can initiate, and either can be the first to commit.
\end{intuitionbox}

\subsection{How SAE works}

SAE runs in a cyclic group $\G$ of prime order $q$ -- a large prime-field group or,
in practice, an elliptic curve like NIST P-256 -- with generator $g$. Both sides
share the password $\pi$. The exchange has two phases: \emph{commit} and
\emph{confirm}.

\textbf{Setup -- hashing the password to the group.} First both sides derive a
group element from the password, the ``password element'' $P = \mathrm{H2C}(\pi,
\ldots)$ (hash-to-curve / hash-to-group), mixing in the two MAC addresses so $P$ is
specific to this pair of devices. The key property: $P$ is a uniform-looking
group element, and recovering $\pi$ from $P$ is as hard as the discrete log -- you
cannot ``read off'' the password from $P$.

\textbf{Commit phase.} Each party (say Alice) picks two random scalars $r_A,
m_A \in \Z_q$, sets her private scalar $s_A = (r_A + m_A) \bmod q$, and computes
the ``element''
\[
  E_A = -\,m_A \cdot P \qquad (\text{in the group, } P^{-m_A}\text{ multiplicatively}),
\]
then sends the pair $(s_A,\ E_A)$. Bob does the same, sending $(s_B,\ E_B)$.

\textbf{Deriving the shared secret.} Alice computes
\[
  K = r_A \cdot \bigl(s_B \cdot P + E_B\bigr).
\]
Substituting $s_B \cdot P + E_B = (r_B + m_B)P - m_B P = r_B \cdot P$, this is
$K = r_A \cdot r_B \cdot P$. Bob, symmetrically, computes $K = r_B \cdot r_A \cdot
P$ -- the \emph{same} element. That shared $K = r_A r_B P$ is a Diffie-Hellman
secret in the group generated by $P$: each side contributed a fresh random scalar
($r_A$, $r_B$), giving forward secrecy, and $P$ ties the whole thing to the
password.

\textbf{Confirm phase.} Each side hashes $K$ (plus the transcript) into a key
$\kappa$ and sends a confirmation MAC. If both used the same password, the same
$P$, hence the same $K$, the confirmations verify and the exchange succeeds,
producing the PMK for the (unchanged) four-way handshake that follows. If the
passwords differ, $P$ differs, $K$ differs, and confirmation fails -- with no
information leaked about \emph{which} password was wrong.

\begin{figure}[H]
\centering
\begin{tikzpicture}[font=\scriptsize, >=stealth]
  % Parties
  \node[draw=aliceblue!80, fill=aliceblue!15, rounded corners=2pt, minimum width=3.0cm, minimum height=0.7cm, align=center] (sta) at (0,0) {
    \textbf{Station (Alice)}\\
    {\tiny Knows password $\pi \implies P = \operatorname{pwd2pt}(\pi)$}
  };
  \node[draw=bobgreen!80, fill=bobgreen!15, rounded corners=2pt, minimum width=3.0cm, minimum height=0.7cm, align=center] (ap) at (8.5,0) {
    \textbf{Access Point (Bob)}\\
    {\tiny Knows password $\pi \implies P = \operatorname{pwd2pt}(\pi)$}
  };

  \draw[thick, gray!60] (sta) -- (0, -6.0);
  \draw[thick, gray!60] (ap) -- (8.5, -6.0);

  % Commit phase
  \node[draw=gray!50, fill=gray!10, rounded corners=1pt, font=\tiny, align=left, anchor=west] at (0.2, -1.0) {
    Pick random scalars $r_A, m_A$.\\
    Mask: $s_A = (r_A + m_A) \bmod q$\\
    Element: $E_A = -m_A P$
  };
  \node[draw=gray!50, fill=gray!10, rounded corners=1pt, font=\tiny, align=left, anchor=east] at (8.3, -1.0) {
    Pick random scalars $r_B, m_B$.\\
    Mask: $s_B = (r_B + m_B) \bmod q$\\
    Element: $E_B = -m_B P$
  };

  \draw[->, thick, aliceblue!80!black] (0, -2.1) -- node[above, font=\tiny]{\textbf{Commit}: Scalar $s_A$, Element $E_A$} (8.5, -2.1);
  \draw[<-, thick, bobgreen!80!black] (0, -2.7) -- node[above, font=\tiny]{\textbf{Commit}: Scalar $s_B$, Element $E_B$} (8.5, -2.7);

  % Shared Key derivation
  \node[draw=structurepurple!80, fill=structurepurple!10, rounded corners=2pt, font=\tiny, align=center, text width=3.8cm, anchor=west] at (0.2, -3.7) {
    Compute Shared Secret:\\
    $K = r_A (s_B P + E_B) = r_A r_B P$\\
    Derive keys: $\kappa = \operatorname{KDF}(K)$
  };
  \node[draw=structurepurple!80, fill=structurepurple!10, rounded corners=2pt, font=\tiny, align=center, text width=3.8cm, anchor=east] at (8.3, -3.7) {
    Compute Shared Secret:\\
    $K = r_B (s_A P + E_A) = r_A r_B P$\\
    Derive keys: $\kappa = \operatorname{KDF}(K)$
  };

  % Confirm phase
  \draw[->, thick, aliceblue!80!black] (0, -4.7) -- node[above, font=\tiny]{\textbf{Confirm}: $\operatorname{MAC}_\kappa(\text{Send-Confirm } 1, s_A, E_A, s_B, E_B)$} (8.5, -4.7);
  \draw[<-, thick, bobgreen!80!black] (0, -5.3) -- node[above, font=\tiny]{\textbf{Confirm}: $\operatorname{MAC}_\kappa(\text{Send-Confirm } 2, s_B, E_B, s_A, E_A)$} (8.5, -5.3);

  % Security outcome banner
  \node[draw=targetgreen!70, fill=targetgreen!10, rounded corners=2pt, font=\tiny, align=center, text width=12.0cm] at (4.25, -6.4) {
    \textbf{Zero-Knowledge Property}: Eavesdropper Eve captures $(s_A, E_A)$ and $(s_B, E_B)$. Because $s_A P + E_A = r_A P$\\
    is an ephemeral DH public key, Eve cannot test password guesses offline. Wrong passwords fail the confirm MAC in real-time.
  };
\end{tikzpicture}
\caption{The WPA3 Simultaneous Authentication of Equals (SAE / Dragonfly) protocol. Both parties hash
the passphrase to a curve point $P$ and perform a blinded Diffie--Hellman exchange. Offline dictionary
attacks are mathematically impossible because the public wire transcripts contain no verifier or MIC for password testing.}
\label{ch26:fig-sae}
\end{figure}

\begin{whybox}
Why can't Eve grind offline? WPA2's handshake also mixes in fresh nonces, and Eve
cracks it anyway, so ``fresh randomness'' alone is not the answer. The answer is a
two-step trap (written multiplicatively, as in Example~\ref{ch26:toy-sae}).
\textbf{First, the Commit messages fit every password.} For any guessed element $P'$,
each observed $E_A$ equals $P'^{-m'_A}$ for some exponent $m'_A$, which in turn defines
$r'_A = s_A - m'_A$. So every guess is perfectly consistent with $(s_A, E_A, s_B, E_B)$,
and nothing is ruled out. \textbf{Second, checking a guess against Confirm is a
Diffie--Hellman problem.} With her guess, Eve can unblind the commits into
$U_A = P'^{s_A} E_A = P'^{r'_A}$ and $U_B = P'^{s_B} E_B = P'^{r'_B}$. To test the
captured Confirm MAC she needs $K' = P'^{r'_A r'_B}$ --- computing CDH on
$(P', U_A, U_B)$ for \emph{every single} candidate password. So a password guess can only be tested by \emph{running} the
protocol against a live party who knows the real $P$: one guess per online
attempt, rate-limitable. This is what a PAKE is designed to guarantee, and it is why WPA3
defeats the capture-and-crack attack that dooms weak WPA2 passphrases.
\end{whybox}

\begin{example}
\label{ch26:toy-sae}
A toy SAE in $\Zns{23}$, a cyclic group of order $q = 22$ with generator $g = 5$.
Suppose the password element is $P = g^{7} \bmod 23$. Compute $5^7 = 78125$; since
$5^{22} \equiv 1$, reduce: $5^2 = 25 \equiv 2$, $5^4 \equiv 4$, $5^7 = 5^4 \cdot
5^2 \cdot 5 \equiv 4 \cdot 2 \cdot 5 = 40 \equiv 17 \pmod{23}$. So $P = 17$.

Alice picks $r_A = 3$, $m_A = 5$, so $s_A = 8$ and $E_A = P^{-m_A} = 17^{-5}$.
Bob picks $r_B = 4$, $m_B = 6$, so $s_B = 10$ and $E_B = 17^{-6}$.
As numbers on the wire: $17^5 \equiv 21$ and $21 \cdot 11 = 231 \equiv 1$, so $E_A = 11$;
$17^6 \equiv 12$ and $12 \cdot 2 = 24 \equiv 1$, so $E_B = 2$.

Alice's shared secret (writing the group multiplicatively):
\[
  K = \bigl(P^{s_B} \cdot E_B\bigr)^{r_A}
    = \bigl(P^{r_B + m_B} \cdot P^{-m_B}\bigr)^{r_A}
    = P^{r_B \, r_A} = 17^{12} \bmod 23.
\]
Compute $17 \equiv -6$, so $17^2 \equiv 36 \equiv 13$; $17^4 \equiv 13^2 = 169
\equiv 8$; $17^8 \equiv 64 \equiv 18$; thus $17^{12} = 17^8 \cdot 17^4 \equiv 18
\cdot 8 = 144 \equiv 6 \pmod{23}$. So $K = 6$. Check it from the actual wire values: Alice receives $(s_B, E_B) = (10, 2)$, computes $17^{10} \equiv 4$, and $(4 \cdot 2)^3 = 512 \equiv 6$. Bob computes $K = P^{r_A r_B} =
17^{12} = 6$ too. They agree on $K = 6$ without either sending a password-derived
value an eavesdropper could test.
\end{example}

\begin{trapbox}
Example~\ref{ch26:toy-sae} uses a $23$-element group so the arithmetic fits on a
page; the discrete log there is trivial to brute-force. Real SAE uses P-256 or a
$3072$-bit prime field, where discrete log is infeasible. Also, real hash-to-curve
(deriving $P$ from the password) is the subtle part -- and is exactly where the
first WPA3 attack landed, next.
\end{trapbox}

\subsection{Dragonblood: early attacks on WPA3}

\begin{attackbox}
\textbf{Dragonblood (Vanhoef and Ronen, 2019).} WPA3's first year brought side-
channel attacks on the hash-to-group step. The original SAE used a ``hunting and
pecking'' loop to map the password to a curve point: try $\mathrm{H}(\pi,
\text{counter})$, check if it is a valid x-coordinate, increment the counter if
not. The number of iterations, and the cache-access pattern inside, \emph{depended
on the password}. By measuring timing or cache behavior of an AP, an attacker
learned how many iterations a given password needed -- a partial fingerprint --
and could prune the password space, restoring a form of offline attack. There were
also downgrade attacks forcing WPA3-capable devices back to WPA2 (and thus back to
the crackable handshake), and denial-of-service via the computationally heavy SAE
commit.

The fixes: a constant-time hash-to-curve (the ``hash-to-element'' / SSWU method,
borrowing from the broader hash-to-curve standardization effort) that runs in
fixed time regardless of password, plus transition-mode hardening against
downgrade. WPA3 with these fixes -- ``WPA3 with H2E'' -- is the current
recommendation.
\end{attackbox}

The Dragonblood story carries the same lesson as KRACK: the core cryptographic idea
(a PAKE for WiFi) was sound, and the flaw lay in the \emph{implementation around
it} -- a password-dependent loop that leaked through a side channel. In
cryptography, constant-time code is a correctness requirement.

\section{The practical attacker: deauth, evil twins, and what's actually at risk}

Now consider the operational question: you are sitting in a
cafe on \texttt{CoffeeHouse}-WiFi (WPA2-Personal). Who can do what?

\begin{focusbox}
\textbf{What an attacker within radio range (not on the network) can do.}
\begin{itemize}
  \item \textbf{See all metadata.} Frame headers -- your device's MAC address (
        unless MAC randomization is on), the AP you joined, the sizes and timing
        of your frames, which networks your device probes for -- are all in the
        clear. Traffic analysis is wide open.
  \item \textbf{Deauthenticate you.} In plain WPA2, management frames are
        unauthenticated, so a spoofed deauth knocks you off at will (denial of
        service, or to force a handshake for cracking, or to push you toward an
        evil twin). \textbf{802.11w / Protected Management Frames (PMF)} fixes
        this by MAC-ing management frames; PMF is \emph{mandatory} in WPA3 and
        optional in WPA2. With PMF on, deauth spoofing stops working.
  \item \textbf{Capture a handshake and crack a weak passphrase} offline, as in
        the attack box above. Strong passphrase or WPA3: no.
\end{itemize}

\textbf{What an attacker on the same network (knows the passphrase) can do.}
\begin{itemize}
  \item In \textbf{WPA2-Personal}, everyone shares the same PMK. But each device's
        PTK is fresh (different nonces), so your neighbor cannot \emph{directly}
        decrypt your unicast frames from a passive capture -- almost. The catch:
        if they capture your four-way handshake (its nonces are public) and they
        know the PMK (they do -- same passphrase!), they can derive \emph{your}
        PTK and decrypt your traffic. So on WPA2-Personal, \textbf{anyone who knows
        the WiFi password and can capture your handshake can read your link-layer
        traffic.} The cafe's other customers are not isolated from you. (Forcing a
        deauth gives them your handshake on demand.)
  \item They can run ARP spoofing / become a man-in-the-middle for your local
        traffic, redirect DNS, etc. -- standard LAN attacks.
  \item \textbf{WPA3} and WPA2-Enterprise close this: SAE gives each device an
        independent PMK that other password-holders cannot derive, so same-network
        users \emph{cannot} decrypt each other's traffic.
\end{itemize}
\end{focusbox}

\begin{attackbox}
\textbf{The evil twin.} Eve sets up her own AP broadcasting the SSID
\texttt{CoffeeHouse} -- same name, stronger signal (she is closer, or boosts
power). She deauths Alice from the real AP. Alice's device, seeing a familiar
network name with a strong signal, may reconnect to Eve's AP. If Eve's AP is open
(no password) or Eve knows/guessed the passphrase, Alice now routes all her
traffic through Eve.

What saves Alice: WPA2/WPA3 \emph{mutual} authentication. In the four-way
handshake, the AP must prove it knows the PMK (message M3's MIC). Eve's evil twin,
\emph{without the passphrase}, cannot produce a valid M3, so a WPA2 client will
not complete the handshake -- the evil twin fails against a password-protected
network it cannot authenticate to. This is why evil-twin attacks target
\textbf{open} networks, captive portals, or networks whose passphrase the attacker
already knows, and why they often pair with phishing (``enter the WiFi password to
continue'') to get the missing secret. WPA3's mandatory PMF and SAE raise the bar
further by killing the deauth that bootstraps the attack.
\end{attackbox}

\begin{intuitionbox}
The reassuring takeaway: even on a hostile WiFi network, your \emph{application}
security can hold. WiFi protects the link; TLS protects the connection end to end.
When you load \texttt{https://bank.com}, the TLS handshake authenticates the bank
with a certificate that Eve's evil twin cannot forge, and encrypts everything
above the link layer. An attacker who fully owns the WiFi sees TLS ciphertext and
SNI/metadata -- annoying, but not your password or your balance. This is the
layered-defense principle: WiFi crypto and TLS guard different layers, and you
want both. The danger zone is unencrypted application traffic (plain HTTP, some
IoT protocols, DNS without DoH) on a hostile link -- there, link-layer compromise
is full compromise.
\end{intuitionbox}

\begin{applicationbox}
\textbf{In deployment, what to actually do.}
\begin{itemize}
  \item \textbf{Home / small office.} Use WPA3-Personal if all your devices
        support it; otherwise WPA2/WPA3 ``transition mode'' with a \emph{long
        random} passphrase (the passphrase is your real security on WPA2). Enable
        PMF. Never WEP or WPA (TKIP); disable them on the router.
  \item \textbf{Enterprise / campus.} WPA2/WPA3-Enterprise with 802.1X/EAP-TLS or
        PEAP against RADIUS. Each user authenticates individually, gets an
        independent PMK, and there is no shared passphrase to crack or to let
        users decrypt each other. Validate the RADIUS server certificate on
        clients, or evil-twin RADIUS impersonation re-opens the door.
  \item \textbf{Public WiFi (you, the user).} Assume the link is hostile. Rely on
        TLS/HTTPS everywhere, prefer a VPN, keep your OS patched (KRACK and
        Dragonblood fixes ship in updates), and enable MAC randomization to blunt
        tracking. ``WPA2 with a password'' on a cafe network does \emph{not} hide
        your traffic from other patrons who know that same password.
  \item \textbf{Open networks (Opportunistic Wireless Encryption, OWE).} WPA3
        adds OWE / ``Enhanced Open,'' which runs an unauthenticated Diffie-Hellman
        on open networks so at least passive eavesdropping is blocked even with no
        password. It does not stop an active evil twin, but it means open
        coffee-shop WiFi no longer has to be plaintext.
\end{itemize}
\end{applicationbox}

\section{Putting the layers together}

Taken together, the stack has several layers. At the root is a
password, low entropy and human-chosen -- the weakest link, and the one every
attack aims at. WPA2 stretches it (PBKDF2) and proves knowledge of it (four-way
handshake MICs), but cannot escape that a captured handshake makes the password
grindable offline. WPA3 breaks that chain with a PAKE (SAE), so the password is no
longer offline-testable and each session gets a forward-secret key. On top of
either sits the data-frame protection -- AES-CCMP (WPA2) or AES-GCMP (WPA3) --
whose entire safety rests on a never-repeating packet number, the invariant KRACK
showed is easy to violate by accident.

\begin{center}
\begin{tikzpicture}[font=\small, node distance=0.35cm,
  layer/.style={draw,thick,rounded corners,minimum width=10.5cm,
                minimum height=0.95cm,align=center}]
  \node[layer,fill=red!8]    (pw)  {Passphrase (low entropy) -- the attacker's target};
  \node[layer,fill=orange!10,below=of pw] (pmk)
       {WPA2: PBKDF2 $\to$ shared PMK \quad / \quad WPA3: SAE PAKE $\to$ per-session PMK (forward-secret)};
  \node[layer,fill=yellow!12,below=of pmk] (hs)
       {Four-way handshake: nonces $\to$ PTK (KCK $\|$ KEK $\|$ TK), mutual MIC proof};
  \node[layer,fill=green!12,below=of hs] (data)
       {Data frames: AES-CCMP (WPA2) / AES-GCMP-256 (WPA3), packet number = nonce};
  \node[layer,fill=blue!10,below=of data] (mgmt)
       {Management frames: PMF (802.11w) -- mandatory in WPA3, blocks deauth spoofing};

  \draw[->,>=Stealth,thick] (pw) -- (pmk);
  \draw[->,>=Stealth,thick] (pmk) -- (hs);
  \draw[->,>=Stealth,thick] (hs) -- (data);
\end{tikzpicture}
\end{center}

Every famous WiFi break maps onto one of these layers. WEP failed at the data
layer (keystream reuse, linear checksum) and the key layer (FMS). The offline
dictionary attack lives at the password layer and is intrinsic to WPA2-Personal.
KRACK struck the handshake/data boundary by resetting a nonce. Dragonblood hit
WPA3's password-to-group mapping through a side channel. The protocols got better
by hardening each layer in turn. The same lesson appears at every layer: a sound
cryptographic core can still be undone by the engineering around it -- a
checksum that should be a MAC, a counter that resets, a loop whose length leaks the
secret. Cryptographic engineering means getting all of these details right at once,
even on a $\$30$ router.

WiFi protects the radio hop to the router and TLS protects the connection to a server,
yet the company running that server still sees every plaintext message it relays. In
Chapter~27, we turn to end-to-end encrypted messaging and the Signal protocol, where X3DH
and the Double Ratchet keep messages secret even from the server and when recipients are
offline.

\section{Exercises}

\begin{exercise}
\textbf{(WEP keystream reuse.)} Two WEP frames are captured with the same IV. Their
ciphertexts are $c_1$ and $c_2$. You know that $m_1$ is an ARP request whose first
$8$ bytes are the constant \texttt{AA AA 03 00 00 00 08 06}. Show exactly how to
recover the first $8$ bytes of $m_2$, and explain why no part of the WEP key is
needed.
\end{exercise}

\begin{exercise}
\textbf{(CRC malleability.)} WEP's integrity check is CRC-32, which is linear:
$\mathrm{CRC}(x \xor \delta) = \mathrm{CRC}(x) \xor \mathrm{CRC}(\delta)$.
Encryption is $c = (m \concat \mathrm{CRC}(m)) \xor Z$. Given a ciphertext $c$ for
an unknown $m$, describe how to produce a valid ciphertext $c'$ that decrypts to
$m \xor \delta$ for a chosen $\delta$, \emph{without} knowing $m$ or the keystream
$Z$. Why does this break integrity?
\end{exercise}

\begin{exercise}
\textbf{(PBKDF2 cost.)} A GPU rig tests $2 \times 10^{6}$ WPA2 passphrase
candidates per second. (a) How long to exhaust all $8$-character lowercase-letter
passphrases ($26^8$)? (b) All $8$-character passphrases over a $72$-character
alphabet ($72^8$)? (c) A five-word diceware passphrase ($7776^5$ possibilities)?
Comment on which are realistically crackable.
\end{exercise}

\begin{exercise}
\textbf{(Toy handshake.)} Using the toy PRF $\mathrm{prf}(K,x) = (Kx + 7) \bmod
101$ from Example~\ref{ch26:toy-handshake}, let $\text{PMK} = 30$, $\text{Anonce}
= 20$, $\text{Snonce} = 15$ (combine by addition). Compute the toy PTK, then
$\text{KCK} = \text{PTK} \bmod 10$, then the M2 MIC $\mathrm{prf}(\text{KCK},
\text{Snonce})$. Now show that, because this toy PRF is linear, an attacker who
sees one (input, output) pair under KCK can solve for KCK -- and explain in one
sentence why HMAC-SHA1 does not have this defect.
\end{exercise}

\begin{exercise}
\textbf{(KRACK invariant.)} State the single invariant whose violation enables
KRACK. Then explain why the official patch -- ``do not reset the packet number on
key reinstall'' -- restores it, and why this required no change to the key-
derivation math or the cipher.
\end{exercise}

\begin{exercise}
\textbf{(Nonce reuse, two modes.)} CCMP uses AES-CCM; WPA3's high tier uses AES-
GCM. Reusing a nonce is bad in both, but worse in one. Explain what an attacker
gains from nonce reuse in CTR-based confidentiality (both modes), and what
\emph{additional} catastrophe GCM suffers that CCM does not. (Hint: GCM's
authenticator is a polynomial evaluation whose key the nonce reuse can expose.)
\end{exercise}

\begin{exercise}
\textbf{(SAE arithmetic.)} Work a toy SAE in $\Zns{23}$ ($g = 5$, $q = 22$) with
password element $P = g^{9} \bmod 23$. Let Alice use $r_A = 5, m_A = 4$ and Bob
use $r_B = 6, m_B = 3$. Compute $s_A, s_B$, the shared secret $K = P^{r_A r_B}
\bmod 23$, and verify both sides reach the same $K$. Then argue informally why an
eavesdropper who sees $s_A, s_B, E_A, E_B$ cannot recover $P$ (and hence the
password) when the group is large.
\end{exercise}

\begin{exercise}
\textbf{(Evil twin reasoning.)} Explain precisely why a WPA2-Personal evil twin
\emph{without} the passphrase cannot complete the four-way handshake with a victim
client, identifying which message it fails to produce and why. Then explain why an
\emph{open} (passwordless) network gives a victim no such protection, and what OWE
(Enhanced Open) does and does not fix.
\end{exercise}

\begin{exercise}
\textbf{(Same-network decryption.)} On a WPA2-Personal cafe network, you and a
stranger both know the passphrase. The stranger captures your four-way handshake.
Walk through exactly what they can derive (PMK, then PTK, then TK) and conclude
whether they can decrypt your unicast traffic. Then explain why the same attack
fails on WPA3-Personal even though, in a sense, ``everyone still knows the
password.''
\end{exercise}

\begin{exercise}
\textbf{(Forward secrecy.)} Define forward secrecy. Explain why WPA2-Personal lacks
it (what an attacker who records traffic today and learns the passphrase next year
can do), and why WPA3's SAE provides it. Identify the specific ephemeral values in
SAE that make recorded sessions undecryptable even after password compromise.
\end{exercise}

\begin{exercise}
\textbf{(Layered defense.)} You connect to a fully attacker-controlled WiFi AP
(the attacker owns the link entirely). For each of the following, state whether the
attacker can read or tamper with the content, and why: (a) an \texttt{https://}
page load to a site with a valid certificate; (b) a plain \texttt{http://} page
load; (c) a DNS lookup over plain UDP; (d) a DNS lookup over DoH (DNS-over-HTTPS).
Use the result to justify the advice ``assume public WiFi is hostile and rely on
TLS.''
\end{exercise}
